\documentclass[12pt]{article}
\usepackage[utf8]{inputenc}
\usepackage[T1]{fontenc}
\usepackage{hyperref}
\title{Tobian Language Resources for K-12 Curriculum Development}
\author{Compiled by Claude Sonnet 5 (claude-sonnet-5) for Andr\'as Kornai}
\date{}
\begin{document}
\maketitle

\section{Introduction}
Tobian (Ramari Hatohobei; ISO 639-3 tox) is a Micronesian language of Hatohobei state, Palau. Speaker estimates in the search results vary, from about 22 native speakers to about 150 people. The ELAR archive describes the language as severely endangered.

\section{Primary Educational Resources}
No school textbooks or primers were found in the search results. The Friends of Tobi pages below may hold teaching-relevant word lists, but I did not verify this.
\begin{itemize}
\item Tobian Language Project (Friends of Tobi). \url{https://friendsoftobi.org/wordweek/tobianlanguageproject.html}
\item Word of the Week research report, 1 March 2015. \url{https://tobi.gmu.edu/wordweek/researchreport03012015.pdf}
\end{itemize}

\section{Linguistic and Reference Materials}
\begin{itemize}
\item ``Toward a Tobian-English Dictionary'', compiled by Barbara Webster Black. It appears to be a work in progress. \url{https://www.friendsoftobi.org/TobianLexicon/Dictionary.pdf}
\item Wikipedia article on the Tobian language. \url{https://en.wikipedia.org/wiki/Tobian_language}
\item Rosetta Panglossia entry for ISO code tox. \url{https://rosettapanglossia.longnow.org/index.php?title=Tox}
\end{itemize}
Older sources mentioned in the search summary, without URLs: Capell (1951, CIMA Report 30) and van den Berg (2014), \emph{Linguistic Sketch of Hatohobeian and Sonsorolese}. These are not cited here as resources.

\section{Religious and Community Resources}
No scripture or religious materials were found in the search results. The Friends of Tobi website is the main community-linked source seen.

\section{Digital and Audio Resources}
\begin{itemize}
\item ELAR (SOAS) deposit ``Documenting Ramari Hatohobei, the Tobian language, a severely endangered Micronesian language'', with video, photographs and documents. \url{https://eldp.access.preservica.com/dk0295/}
\item Language in Use: documentation work on Tobian by Peter and Isauro. \url{https://tobi.gmu.edu/wordweek/languageinuse.htm}
\end{itemize}

\section{Curriculum Development Context}
The language has very few speakers, so a curriculum depends on the documentation archive and the draft dictionary. The Friends of Tobi pages say Justin Andrew has agreed to maintain the dictionary and the archive deposit. No school programme was found.
\end{document}